\documentclass[12pt,a4paper,openany,oneside]{book}

% --- Paquetes ---
\usepackage[utf8]{inputenc}
\usepackage[spanish, es-noshorthands]{babel}
\usepackage{amsmath, amsfonts, amssymb}
\usepackage{graphicx}
\usepackage{geometry}
\usepackage{hyperref}
\usepackage{booktabs}
\usepackage{fancyhdr}
\usepackage{listings}
\usepackage{xcolor}
\usepackage{tikz}
\usepackage{pgfplots}
\usepackage{colortbl}
\usepackage{multirow}
\usepackage{hhline}
\usetikzlibrary{arrows.meta, positioning, shapes.geometric, calc, shapes}
\pgfplotsset{compat=1.18}

\geometry{margin=2.5cm}

% --- Colores y estilos para código Python ---
\definecolor{codegreen}{rgb}{0,0.6,0}
\definecolor{codegray}{rgb}{0.5,0.5,0.5}
\definecolor{codepurple}{rgb}{0.58,0,0.82}
\definecolor{backcolour}{rgb}{0.95,0.95,0.92}

\lstset{
    language=Python,
    backgroundcolor=\color{backcolour},   
    commentstyle=\color{codegreen},
    keywordstyle=\color{blue},
    numberstyle=\tiny\color{codegray},
    stringstyle=\color{codepurple},
    basicstyle=\ttfamily\small,
    breaklines=true,
    frame=single,
    showstringspaces=false
}

% --- Estilo de cabecera y pie ---
\pagestyle{fancy}
\fancyhf{}
\renewcommand{\headrulewidth}{0.4pt}
\renewcommand{\footrulewidth}{0.4pt}
\fancyhead[L]{\small Carlos César Sánchez Coronel}
\fancyhead[R]{\small Sesión 6: Redes Neuronales para NLP}
\fancyfoot[L]{\small Carlos César Sánchez Coronel}
\fancyfoot[C]{\thepage}

% --- Portada ---
\title{
    \textbf{Sesión 6} \\ 
    \Huge \textbf{REDES NEURONALES PARA NLP Y EMBEDDINGS ENTRENABLES} \\
    \Large De representaciones estáticas a aprendizaje end-to-end
}
\author{
    \small Por \\ \vspace{0.2cm}
    \Large \textsc{Carlos César Sánchez Coronel}
}
\date{2026}

\begin{document}

\maketitle
\tableofcontents
\addcontentsline{toc}{chapter}{Índice general}

% ------------------------------------------------------------------
% CAPÍTULO 6: SESIÓN 6 - REDES NEURONALES PARA NLP
% ------------------------------------------------------------------
\chapter{Redes Neuronales para NLP y Embeddings Entrenables}

En las sesiones anteriores se han utilizado representaciones de texto (BOW, TF-IDF, word embeddings estáticos) junto con modelos clásicos de machine learning. Sin embargo, estos enfoques tienen una limitación fundamental: la representación de las palabras es fija y no se ajusta a la tarea específica. Esta sesión introduce las redes neuronales para NLP, comenzando por la capa de embedding entrenable y las arquitecturas feed-forward simples. Se sientan las bases para modelos más complejos como RNNs y Transformers, que se abordarán en sesiones posteriores.

\section{Logro de la sesión}
Construir redes neuronales simples para tareas de NLP usando PyTorch, comprendiendo el concepto de embeddings entrenables, el flujo de datos desde el texto hasta la predicción, y las técnicas de regularización específicas para NLP.

\section{Introducción: de los embeddings estáticos a los entrenables}

\subsection{Recordatorio: embeddings estáticos}
Word2Vec y GloVe producen embeddings estáticos: una vez entrenados, el vector para una palabra es fijo, independientemente del contexto o de la tarea posterior. Estos embeddings se pueden usar como características de entrada para modelos clásicos (SVM, etc.) o para inicializar capas de embedding en redes neuronales.

\subsection{Limitaciones de los embeddings estáticos}
\begin{itemize}
    \item No se ajustan a la tarea específica (ej. la palabra "banco" puede necesitar un significado diferente en un clasificador de textos financieros vs. uno de biología).
    \item No capturan información contextual más allá de la ventana de entrenamiento original.
    \item Son independientes del dominio; un embedding entrenado en Wikipedia puede no ser óptimo para tweets médicos.
\end{itemize}

\subsection{Ventaja de los embeddings entrenables}
En una red neuronal, la capa de embedding puede ser \textbf{entrenable}: sus pesos se actualizan durante el backpropagation para minimizar la pérdida de la tarea final. Esto permite que el modelo aprenda representaciones específicas para el problema, mejorando el rendimiento.

\section{Capa de Embedding en PyTorch}

\subsection{Fundamento matemático}
Una capa de embedding es esencialmente una tabla de búsqueda. Para un vocabulario de tamaño \(|V|\) y una dimensión de embedding \(d\), se define una matriz \(\mathbf{E} \in \mathbb{R}^{|V| \times d}\). Cada fila \(\mathbf{E}[i]\) corresponde al embedding de la palabra con índice \(i\).

Dado un tensor de índices de palabras \(\mathbf{x} \in \mathbb{Z}^{n}\) (donde \(n\) es la longitud de la secuencia), la capa de embedding produce un tensor \(\mathbf{X} \in \mathbb{R}^{n \times d}\) tal que:
\[
\mathbf{X}[j] = \mathbf{E}[ \mathbf{x}[j] ]
\]

\subsection{Implementación en PyTorch}
\begin{lstlisting}[language=Python]
import torch
import torch.nn as nn

embedding_dim = 100
vocab_size = 10000
embedding_layer = nn.Embedding(vocab_size, embedding_dim)

# Ejemplo: secuencia de 5 palabras con índices
input_indices = torch.tensor([1, 5, 23, 7, 42])
embedded = embedding_layer(input_indices)  # shape: (5, 100)
\end{lstlisting}

\subsection{Inicialización}
Por defecto, los embeddings se inicializan aleatoriamente (distribución uniforme o normal). Se pueden inicializar con embeddings pre-entrenados (Word2Vec, GloVe) y luego decidir si se congelan o se ajustan durante el entrenamiento.

\begin{lstlisting}[language=Python]
# Cargar embeddings pre-entrenados (ej. GloVe) y asignarlos
pretrained_weights = torch.tensor(glove_vectors)  # shape: (vocab_size, dim)
embedding_layer = nn.Embedding.from_pretrained(pretrained_weights, freeze=False)
# freeze=True para mantenerlos estáticos
\end{lstlisting}

\section{Redes Feed-Forward para Clasificación de Texto}

\subsection{Arquitectura básica}
Una arquitectura simple pero efectiva para clasificación de texto consta de:
\begin{enumerate}
    \item \textbf{Capa de embedding}: convierte índices de palabras a vectores densos.
    \item \textbf{Función de pooling}: combina los embeddings de todas las palabras de la secuencia en un único vector (promedio, suma, max-pooling).
    \item \textbf{MLP (Multilayer Perceptron)}: una o más capas fully-connected con activaciones no lineales.
    \item \textbf{Capa de salida}: con activación softmax para obtener probabilidades por clase.
\end{enumerate}

\subsection{Promedio de embeddings (Average Pooling)}
La forma más simple de combinar los embeddings de una oración es calcular su promedio:

Dada una secuencia de \(T\) palabras con embeddings \(\mathbf{e}_1, \mathbf{e}_2, ..., \mathbf{e}_T \in \mathbb{R}^d\), el vector de la oración es:
\[
\mathbf{h} = \frac{1}{T} \sum_{t=1}^T \mathbf{e}_t
\]

\textbf{Ventaja}: invariante a la longitud de la oración.
\textbf{Desventaja}: pierde el orden y el contexto local.

\subsection{Arquitectura completa}
\begin{lstlisting}[language=Python]
class TextClassifier(nn.Module):
    def __init__(self, vocab_size, embed_dim, hidden_dim, num_classes, dropout=0.5):
        super().__init__()
        self.embedding = nn.Embedding(vocab_size, embed_dim)
        self.fc1 = nn.Linear(embed_dim, hidden_dim)
        self.fc2 = nn.Linear(hidden_dim, num_classes)
        self.dropout = nn.Dropout(dropout)
        self.relu = nn.ReLU()
        
    def forward(self, x):
        # x: (batch_size, seq_len)
        embedded = self.embedding(x)          # (batch_size, seq_len, embed_dim)
        pooled = embedded.mean(dim=1)         # (batch_size, embed_dim)
        h = self.dropout(self.relu(self.fc1(pooled)))  # (batch_size, hidden_dim)
        out = self.fc2(h)                     # (batch_size, num_classes)
        return out
\end{lstlisting}

\subsection{Consideraciones sobre longitudes de secuencia}
Las oraciones tienen longitudes variables. En PyTorch, se pueden:
\begin{itemize}
    \item Agrupar por longitud y usar `pack_padded_sequence` (más eficiente).
    \item Truncar o rellenar (padding) todas las secuencias a la misma longitud.
    \item Usar máscaras para ignorar el padding en el promedio.
\end{itemize}

\begin{lstlisting}[language=Python]
# Ejemplo con padding y máscara
def masked_mean(embeddings, mask):
    # embeddings: (batch, seq_len, dim)
    # mask: (batch, seq_len) con 1 para tokens reales, 0 para padding
    mask = mask.unsqueeze(-1)  # (batch, seq_len, 1)
    sum_embeddings = (embeddings * mask).sum(dim=1)
    lengths = mask.sum(dim=1)  # (batch, 1)
    return sum_embeddings / (lengths + 1e-8)
\end{lstlisting}

\section{Función de Pérdida: Cross-Entropy}

\subsection{Definición matemática}
Para un problema de clasificación con \(C\) clases, la entropía cruzada para una muestra se define como:

\[
\mathcal{L} = -\sum_{c=1}^C y_c \log(\hat{y}_c)
\]
donde \(y_c\) es 1 para la clase verdadera y 0 en otro caso (one-hot), y \(\hat{y}_c\) es la probabilidad predicha por el modelo (después de softmax).

Para un lote (batch) de \(N\) muestras, la pérdida promedio es:

\[
\mathcal{L}_{\text{batch}} = -\frac{1}{N} \sum_{i=1}^N \sum_{c=1}^C y_{ic} \log(\hat{y}_{ic})
\]

\subsection{Implementación en PyTorch}
PyTorch proporciona `nn.CrossEntropyLoss`, que combina log-softmax y entropía cruzada en una sola función numéricamente estable.

\begin{lstlisting}[language=Python]
criterion = nn.CrossEntropyLoss()
logits = model(x)          # (batch_size, num_classes)
loss = criterion(logits, y)  # y: (batch_size) con índices de clase
\end{lstlisting}

\subsection{Interpretación probabilística}
Minimizar la entropía cruzada equivale a maximizar la log-verosimilitud de los datos bajo el modelo.

\section{Regularización en NLP}

Los modelos de NLP tienden a sobreajustar debido a la alta dimensionalidad del vocabulario. Las técnicas de regularización son esenciales.

\subsection{Dropout}
Durante el entrenamiento, se apagan aleatoriamente neuronas con probabilidad \(p\). Esto evita la co-adaptación y fuerza a la red a aprender representaciones más robustas.

\begin{lstlisting}[language=Python]
self.dropout = nn.Dropout(0.5)
h = self.dropout(self.relu(self.fc1(pooled)))
\end{lstlisting}

\textbf{Efecto}: en inferencia, todas las neuronas se usan pero con pesos escalados por \(1-p\).

\subsection{Weight Decay (Regularización L2)}
Se añade un término a la función de pérdida que penaliza pesos grandes:

\[
\mathcal{L}_{\text{total}} = \mathcal{L}_{\text{original}} + \lambda \sum_{w} \|w\|^2
\]

En PyTorch, se especifica en el optimizador:

\begin{lstlisting}[language=Python]
optimizer = torch.optim.Adam(model.parameters(), lr=0.001, weight_decay=1e-5)
\end{lstlisting}

\subsection{Early Stopping}
Se detiene el entrenamiento cuando la métrica en validación deja de mejorar durante un número de épocas (paciencia). Previene sobreajuste por exceso de entrenamiento.

\subsection{Normalización de capas (Layer Normalization)}
Útil en NLP para estabilizar el entrenamiento, especialmente en modelos recurrentes. Normaliza las activaciones a lo largo de la dimensión de características.

\begin{lstlisting}[language=Python]
self.layer_norm = nn.LayerNorm(embed_dim)
\end{lstlisting}

\subsection{Regularización específica para embeddings}
\begin{itemize}
    \item \textbf{Embedding dropout}: se apagan aleatoriamente palabras completas (útil para evitar que el modelo dependa de palabras específicas).
    \item \textbf{Max-norm regularization}: limita la norma de los vectores embedding a un valor máximo.
\end{itemize}

\section{Ejemplo paso a paso: Clasificación de Sentimiento con PyTorch}

\subsection{Datos}
Utilizaremos un pequeño conjunto de reseñas de películas con etiquetas positivas (1) y negativas (0).

\begin{lstlisting}[language=Python]
texts = [
    "I loved this movie, it was fantastic",
    "Terrible film, I hated it",
    "Amazing acting and great story",
    "Boring and slow, waste of time"
]
labels = [1, 0, 1, 0]
\end{lstlisting}

\subsection{Paso 1: Construir vocabulario}
\begin{lstlisting}[language=Python]
from collections import Counter
import torch
import torch.nn as nn
import torch.optim as optim

def build_vocab(texts, max_size=10000):
    word_counts = Counter()
    for text in texts:
        word_counts.update(text.lower().split())
    vocab = {word: idx+2 for idx, (word, _) in enumerate(word_counts.most_common(max_size))}
    vocab['<PAD>'] = 0
    vocab['<UNK>'] = 1
    return vocab

vocab = build_vocab(texts)
vocab_size = len(vocab)
print(f"Tamaño del vocabulario: {vocab_size}")
\end{lstlisting}

\subsection{Paso 2: Tokenizar y convertir a índices}
\begin{lstlisting}[language=Python]
def encode(text, vocab, max_len=10):
    tokens = text.lower().split()
    indices = [vocab.get(token, vocab['<UNK>']) for token in tokens[:max_len]]
    # padding
    indices += [vocab['<PAD>']] * (max_len - len(indices))
    return indices

X = [encode(text, vocab) for text in texts]
X = torch.tensor(X)  # (4, 10)
y = torch.tensor(labels)
\end{lstlisting}

\subsection{Paso 3: Definir el modelo}
\begin{lstlisting}[language=Python]
class SentimentClassifier(nn.Module):
    def __init__(self, vocab_size, embed_dim=50, hidden_dim=32, num_classes=2, dropout=0.5):
        super().__init__()
        self.embedding = nn.Embedding(vocab_size, embed_dim, padding_idx=0)
        self.fc1 = nn.Linear(embed_dim, hidden_dim)
        self.fc2 = nn.Linear(hidden_dim, num_classes)
        self.dropout = nn.Dropout(dropout)
        self.relu = nn.ReLU()
        
    def forward(self, x):
        # x: (batch, seq_len)
        embedded = self.embedding(x)          # (batch, seq_len, embed_dim)
        pooled = embedded.mean(dim=1)         # (batch, embed_dim)
        h = self.dropout(self.relu(self.fc1(pooled)))  # (batch, hidden_dim)
        out = self.fc2(h)                     # (batch, num_classes)
        return out

model = SentimentClassifier(vocab_size)
\end{lstlisting}

\subsection{Paso 4: Entrenamiento}
\begin{lstlisting}[language=Python]
criterion = nn.CrossEntropyLoss()
optimizer = optim.Adam(model.parameters(), lr=0.01, weight_decay=1e-5)

epochs = 100
for epoch in range(epochs):
    optimizer.zero_grad()
    logits = model(X)
    loss = criterion(logits, y)
    loss.backward()
    optimizer.step()
    
    if (epoch+1) % 20 == 0:
        print(f'Epoch {epoch+1}, Loss: {loss.item():.4f}')
\end{lstlisting}

\subsection{Paso 5: Evaluación}
\begin{lstlisting}[language=Python]
with torch.no_grad():
    logits = model(X)
    predictions = torch.argmax(logits, dim=1)
    accuracy = (predictions == y).float().mean()
    print(f'Accuracy: {accuracy:.2%}')
\end{lstlisting}

\section{Métricas de Evaluación en NLP}

\subsection{Accuracy}
Porcentaje de predicciones correctas. Adecuado para clases balanceadas.

\subsection{Precisión, Recall y F1-Score}
Para problemas desbalanceados (ej. detección de spam), son más informativas.

\begin{lstlisting}[language=Python]
from sklearn.metrics import precision_recall_fscore_support

precision, recall, f1, _ = precision_recall_fscore_support(y, predictions, average='binary')
print(f'Precision: {precision:.2f}, Recall: {recall:.2f}, F1: {f1:.2f}')
\end{lstlisting}

\subsection{Matriz de confusión}
Visualiza errores por clase.

\subsection{Curva ROC y AUC}
Para problemas binarios, mide la capacidad discriminativa.

\subsection{Log-Loss (Cross-Entropy)}
La pérdida en sí misma puede ser una métrica, especialmente cuando se necesitan probabilidades bien calibradas.

\section{Aplicaciones de Negocio y Casos de Uso}

\subsection{Análisis de sentimiento en redes sociales}
Las empresas monitorean menciones en Twitter, Facebook, etc., para detectar crisis de reputación o medir aceptación de productos. Una red feed-forward con embeddings entrenables puede superar a modelos clásicos.

\subsection{Clasificación de tickets de soporte}
Automatizar la asignación de tickets a departamentos (ventas, soporte técnico, facturación) basado en el texto del cliente.

\subsection{Detección de spam en correos electrónicos}
Modelos neuronales pueden adaptarse a nuevos patrones de spam mejor que reglas fijas.

\subsection{Moderación de contenido}
Clasificar comentarios en foros como apropiados o inapropiados (tóxicos, ofensivos).

\subsection{Análisis de intenciones en chatbots}
Determinar si un usuario quiere comprar, reclamar, preguntar, etc.

\section{Preparación para Entrevistas}

\subsection{Preguntas típicas}
\begin{itemize}
    \item \textbf{"¿Cómo representarías una oración para ingresarla a una red neuronal?"}
    Respuesta: convertir palabras a índices, usar una capa de embedding (entrenable o pre-entrenada), y luego combinar los embeddings mediante promedio, suma, max-pooling, o usar arquitecturas más complejas como RNNs o Transformers.
    
    \item \textbf{"¿Qué son los embeddings entrenables y por qué son mejores que los estáticos?"}
    Los embeddings entrenables se ajustan durante el entrenamiento de la tarea específica, capturando matices del dominio y mejorando el rendimiento. Los estáticos son fijos y pueden no ser óptimos para la tarea.
    
    \item \textbf{"Explica la función de pérdida cross-entropy."}
    Mide la diferencia entre la distribución de probabilidad predicha y la distribución real (one-hot). Penaliza más las predicciones seguras pero erróneas.
    
    \item \textbf{"¿Cómo evitarías el sobreajuste en un modelo de NLP?"}
    Usar dropout, weight decay, early stopping, aumentar datos (data augmentation), y reducir la complejidad del modelo.
    
    \item \textbf{"¿Qué es el padding y cómo se maneja?"}
    Para procesar oraciones de diferentes longitudes en lotes, se añaden tokens de padding (ej. <PAD>) para que todas tengan la misma longitud. Luego, se usa una máscara para que el modelo ignore esos tokens en operaciones como el promedio o en capas recurrentes.
\end{itemize}

\subsection{Preguntas avanzadas}
\begin{itemize}
    \item \textbf{"¿Cómo inicializarías una capa de embedding con Word2Vec y por qué?"}
    Se cargan los pesos pre-entrenados y se asigna a la capa. Luego se puede decidir congelarlos (si se tiene poco data) o permitir que se ajusten (fine-tuning) si se dispone de suficiente data específica.
    
    \item \textbf{"¿Qué diferencias hay entre usar el promedio de embeddings y usar una LSTM?"}
    El promedio es invariante al orden y pierde información secuencial; es adecuado para tareas donde el orden no importa (ej. clasificación de temas). LSTM captura dependencias temporales, mejor para tareas como análisis de sentimiento donde el orden es crucial.
\end{itemize}

\section{Conexión con el resto del curso}
\begin{itemize}
    \item \textbf{Sesión 4 (Word Embeddings)}: proporciona los embeddings estáticos que pueden usarse como inicialización.
    \item \textbf{Sesión 5 (Modelos Clásicos)}: las redes neuronales superan a estos modelos en rendimiento, pero requieren más datos y cómputo.
    \item \textbf{Sesión 7 (RNN/LSTM)}: extensión natural para capturar orden y contexto.
    \item \textbf{Sesión 8 (Transformers)}: evolución que supera las limitaciones de las RNNs.
\end{itemize}

\section{Resumen}
\begin{itemize}
    \item Los \textbf{embeddings entrenables} se ajustan a la tarea, mejorando el rendimiento respecto a los estáticos.
    \item La \textbf{capa de embedding} es una tabla de búsqueda diferenciable, implementada en PyTorch como `nn.Embedding`.
    \item Una arquitectura simple de \textbf{red feed-forward} promedia los embeddings y los pasa por un MLP.
    \item La función de pérdida estándar es \textbf{Cross-Entropy}, que mide la diferencia entre probabilidades predichas y reales.
    \item La \textbf{regularización} (dropout, weight decay, early stopping) es crucial para evitar sobreajuste.
    \item Este modelo es un baseline efectivo para muchas tareas de NLP y sirve como base para arquitecturas más complejas.
\end{itemize}

\end{document}